\section{切空间与余切空间}
\label{sec:mfd-tangent-cotangent}

球面放在$\mathbb R^3$中时，曲线的速度可以直接在$\mathbb R^3$中求导，得到一个与球面相切的向量。一般流形没有现成的外部空间，速度必须用流形自身的语言描述。上一节已经看到，同一速度在两张图中的分量由坐标变换的Jacobian矩阵换算。本节先在内点处把速度定义为曲线的等价类，组成切空间，再用一族坐标记录处理边界点。随后把对切向量的线性读数组织成余切空间，比较两者在换图时的不同变换规则。

\subsection{切空间}
\label{sec:mfd-tangent-space}

\paragraph{内点处的曲线速度}
设$p$是$M$的内点。经过$p$的光滑曲线$\gamma:(-\varepsilon,\varepsilon)\to M$，$\gamma(0)=p$，在一张包含$p$的图$(U,\varphi)$中有坐标速度$(\varphi\circ\gamma)'(0)\in\mathbb R^n$。两条曲线可能形状完全不同，却在$p$处有相同的坐标速度；就“此刻往哪里、以多快的速度运动”而言，它们没有区别。

\begin{definition}[切向量与切空间]
\label{def:mfd-tangent-space}
设$p$是光滑流形$M$的内点，$(U,\varphi)$是包含$p$的坐标图。对经过$p$的两条光滑曲线$\gamma_1,\gamma_2$，若$(\varphi\circ\gamma_1)'(0)=(\varphi\circ\gamma_2)'(0)$，则称二者在$p$处等价。等价类$[\gamma]$称为$p$处的\term{切向量}（tangent vector），也记作$\gamma'(0)$；全体切向量组成的集合$T_pM$称为$M$在$p$处的\term{切空间}（tangent space）。
\end{definition}

定义中用了一张特定的图，但等价关系与图无关：若在$\varphi$中坐标速度相同，由式\eqref{eq:mfd-curve-velocity-change}，在另一张图$\psi$中二者的坐标速度都等于同一个矩阵乘以同一个向量，仍然相同。

\paragraph{切空间的线性结构}
切向量是等价类，本来没有加法。下面借用坐标记录把$\mathbb R^n$的线性结构搬过来，并说明搬来的结构与所用的图无关。

\begin{theorem}[切空间是$n$维线性空间]
\label{thm:mfd-tangent-linear}
设$p$是$n$维光滑流形$M$的内点，$(U,\varphi)$是包含$p$的坐标图。映射
\[
\Theta_\varphi:T_pM\to\mathbb R^n,\qquad [\gamma]\mapsto(\varphi\circ\gamma)'(0)
\]
是双射。用它把$\mathbb R^n$的加法与数乘搬到$T_pM$上，所得线性结构与$\varphi$的选择无关，$T_pM$是$n$维实线性空间。
\end{theorem}

\begin{proof}
单射直接来自等价关系的定义。为证满射，任取$\bm v\in\mathbb R^n$。$p$是内点，$\varphi(U)$是包含$\varphi(p)$的欧氏开集，所以$t$足够小时$\varphi(p)+t\bm v\in\varphi(U)$，曲线$\gamma(t)=\varphi^{-1}(\varphi(p)+t\bm v)$有定义且光滑，坐标速度为$\bm v$。

若$(V,\psi)$是另一张包含$p$的图，式\eqref{eq:mfd-curve-velocity-change}说明$\Theta_\psi=\bm J\,\Theta_\varphi$，其中$\bm J=\bm J_{\psi\circ\varphi^{-1}}(\varphi(p))$是可逆矩阵。可逆线性映射保持线性组合，所以按$\Theta_\varphi$规定的$a[\gamma_1]+b[\gamma_2]$，在$\Theta_\psi$下的记录是$a\Theta_\psi[\gamma_1]+b\Theta_\psi[\gamma_2]$，两张图规定的运算相同。
\end{proof}

\paragraph{坐标基与变换规则}
$\mathbb R^n$的标准基$\bm e_1,\ldots,\bm e_n$在$\Theta_\varphi^{-1}$下给出$T_pM$的一组基，记为
\[
\left.\frac{\partial}{\partial x_i}\right|_p=\Theta_\varphi^{-1}(\bm e_i),\qquad i=1,\ldots,n,
\]
称为\term{坐标基}（coordinate basis），简记为$\partial_i|_p$。它就是第$i$条坐标线$t\mapsto\varphi^{-1}(\varphi(p)+t\bm e_i)$的速度。任一切向量写成$\bm v=\sum_i v_i\,\partial_i|_p$，系数$(v_1,\ldots,v_n)=\Theta_\varphi(\bm v)$称为它在这张图中的分量。

换一张图$\psi$，坐标记为$\tilde x_1,\ldots,\tilde x_n$，同一切向量的分量变为
\begin{equation}
\tilde v_k=\sum_{j=1}^n\frac{\partial\tilde x_k}{\partial x_j}(\varphi(p))\,v_j,
\qquad\text{即}\qquad
\left.\frac{\partial}{\partial x_j}\right|_p=\sum_{k=1}^n\frac{\partial\tilde x_k}{\partial x_j}(\varphi(p))\left.\frac{\partial}{\partial\tilde x_k}\right|_p,
\label{eq:mfd-tangent-change}
\end{equation}
其中$\partial\tilde x_k/\partial x_j$是坐标变换$\psi\circ\varphi^{-1}$的Jacobian矩阵的元素。左式说明分量按Jacobian矩阵变换；右式说明旧的第$j$个基向量由Jacobian矩阵第$j$列的元素组合新基向量得到。两种写法表达同一件事：记录变了，向量本身不变。

\paragraph{球面上的核对}
回到$p=(0,\tfrac35,\tfrac45)$，北极图坐标$\bm u=(0,3)$。对坐标变换$\tau(\bm u)=\bm u/\lVert\bm u\rVert_2^2$求导，得
\[
\bm J_\tau(\bm u)=\frac{\lVert\bm u\rVert_2^2\,\bm I-2\bm u\bm u^\top}{\lVert\bm u\rVert_2^4},\qquad
\bm J_\tau(0,3)=\frac{9\bm I-2\operatorname{diag}(0,9)}{81}=\operatorname{diag}\Bigl(\frac19,-\frac19\Bigr).
\]
于是北极图中的坐标基向量$\partial_{u_1}|_p$，在南极图中的分量是$(\tfrac19,0)$。直接验证：北极图中的坐标线$t\mapsto(t,3)$在南极图中是$t\mapsto(t,3)/(t^2+9)$，在$t=0$处的导数为$(\tfrac19,0)$，与矩阵计算一致。

借助嵌入也可以看到这个切向量的模样。对式\eqref{eq:mfd-stereographic-north}中$\sigma_N^{-1}$关于$u_1$在$(0,3)$处求导，得$\mathbb R^3$中的速度$(0.2,0,0)$，它与$p$的内积为零，确实与球面相切。不过这只是嵌入在$\mathbb R^3$中的球面才有的另一种记录；切空间的定义本身没有用到$\mathbb R^3$。嵌入子流形的切空间怎样作为环境切空间的子空间出现，见第~\ref{chap:manifold-relations}章。

\paragraph{边界点处的切空间}
在边界点$p$处，坐标速度只能从单侧曲线得到，定义~\ref{def:mfd-tangent-space}不再适用。注意式\eqref{eq:mfd-tangent-change}只用到坐标变换在$\varphi(p)$处的Jacobian矩阵，而边界点处的Jacobian矩阵由光滑延拓确定，并且与延拓的选择无关。于是可以绕开曲线，直接用坐标记录及其换算规则定义切向量。

\begin{definition}[用坐标记录定义的切空间]
\label{def:mfd-tangent-records}
设$p\in M$，$M$可以带边界。考虑所有二元组$(\varphi,\bm v)$，其中$\varphi$是包含$p$的坐标图，$\bm v\in\mathbb R^n$。规定
\[
(\varphi,\bm v)\sim(\psi,\tilde{\bm v})\quad\Longleftrightarrow\quad
\tilde{\bm v}=\bm J_{\psi\circ\varphi^{-1}}(\varphi(p))\,\bm v,
\]
其中边界点处的Jacobian矩阵取坐标变换任一光滑延拓的Jacobian矩阵。等价类组成的集合仍记为$T_pM$。
\end{definition}

链式法则说明$\sim$是等价关系：自反性来自恒等变换，对称性来自$\bm J_{\varphi\circ\psi^{-1}}=\bm J_{\psi\circ\varphi^{-1}}^{-1}$，传递性来自三张图之间坐标变换的复合。固定一张图，$\bm v\mapsto[(\varphi,\bm v)]$是$\mathbb R^n$到$T_pM$的双射，与定理~\ref{thm:mfd-tangent-linear}同样的论证给出与图无关的线性结构，坐标基与变换规则\eqref{eq:mfd-tangent-change}照旧成立。在内点处，$[\gamma]\mapsto[(\varphi,(\varphi\circ\gamma)'(0))]$把两种定义一一对应起来，并保持线性运算，所以两种定义可以等同。

\paragraph{完整切空间与单侧方向}
在边界点处，$T_pM$仍然是完整的$n$维线性空间，其中既有指向内部的方向，也有指向外部的方向。这种分类与图无关。

\begin{lemma}[内向与外向的坐标无关性]
\label{lem:mfd-inward-outward}
设$p\in\partial M$，$\varphi,\psi$是包含$p$的两张坐标图，$\tilde{\bm v}=\bm J_{\psi\circ\varphi^{-1}}(\varphi(p))\bm v$。则$\tilde v_n$与$v_n$同为正、同为负或同为零。
\end{lemma}

\begin{proof}
记$\tau=\psi\circ\varphi^{-1}$，它的最后一个分量$\tau_n$在$\mathbb H^n$上非负，并由定理~\ref{thm:mfd-boundary-invariance}在$\partial\mathbb H^n$上为零。沿边界方向求导，$j<n$时$\partial\tau_n/\partial x_j=0$；从内侧求$x_n$方向的单侧导数，$\partial\tau_n/\partial x_n\geq0$。于是Jacobian矩阵的最后一行是$(0,\ldots,0,c)$，$c\geq0$。矩阵可逆，所以$c>0$，$\tilde v_n=cv_n$。
\end{proof}

据此，$v_n>0$的切向量称为内向的，$v_n<0$的称为外向的，$v_n=0$的与边界相切。从$p$出发进入$M$的单侧光滑曲线$\gamma:[0,\varepsilon)\to M$，速度满足$v_n\geq0$；反过来，每个$v_n\geq0$的$\bm v$都是曲线$t\mapsto\varphi^{-1}(\varphi(p)+t\bm v)$在$t\geq0$时的速度。所以单侧曲线的速度恰好组成闭半空间$\{v_n\geq0\}$。它对取负不封闭，不是线性空间，这正是边界点的切空间不能只用实际可行的运动来定义的原因。

闭圆盘在边界点$(1,0)$处可以具体看到这些方向。用第~\ref{sec:mfd-boundary}小节的极坐标图$(\theta,s)$，$s=1-r$。向内运动的曲线$(1-t,0)$在图中是$(\theta,s)=(0,t)$，速度$(0,1)$，内向；绕圆周的运动速度是$(1,0)$，与边界相切；向量$(0,-1)$外向，它属于$T_pM$，却不是圆盘内任何曲线在$p$处的速度。

\subsection{余切空间}
\label{sec:mfd-cotangent-space}

知道一个点在球面上以什么速度运动之后，常常只关心速度的某一方面，例如高度变化得多快，或者沿某个方向的分量是多少。这样的读数对速度是线性的：速度加倍，读数加倍；两个速度相加，读数相加。对切向量的线性读数构成另一个线性空间。

\begin{definition}[余切空间]
\label{def:mfd-cotangent-space}
$p\in M$处切空间的对偶空间$T_p^*M=(T_pM)^*$，即$T_pM$上全体线性泛函组成的空间，称为$M$在$p$处的\term{余切空间}（cotangent space），其元素称为\term{余切向量}（cotangent vector）或协向量。
\end{definition}

有限维线性空间的对偶空间与原空间维数相同，所以$T_p^*M$也是$n$维的。对偶空间、对偶基与对偶映射的一般讨论可参考\BookChapterRef{chap:linear-representation}{卷二第2章“线性空间、映射与坐标”}。

\paragraph{对偶基}
坐标基$\partial_1|_p,\ldots,\partial_n|_p$有唯一的对偶基，记为$\mathrm dx_1|_p,\ldots,\mathrm dx_n|_p$，满足
\[
\mathrm dx_i|_p\bigl(\partial_j|_p\bigr)=\delta_{i,j},\qquad\text{从而}\qquad
\mathrm dx_i|_p(\bm v)=v_i .
\]
所以$\mathrm dx_i|_p$的作用就是读出切向量在这张图中的第$i$个分量。任一余切向量写成$\omega=\sum_i a_i\,\mathrm dx_i|_p$，系数$a_i=\omega(\partial_i|_p)$是它对第$i$条坐标线速度的读数，作用在切向量上为$\omega(\bm v)=\sum_ia_iv_i$。记号$\mathrm dx_i$的来历在第~\ref{sec:mfd-smooth-functions}小节说明：它正是坐标函数$x_i$的微分。

\paragraph{变换规则的比较}
设另一张图中$\omega=\sum_k\tilde a_k\,\mathrm d\tilde x_k|_p$。把$\omega$作用在$\partial_j|_p$上，并用式\eqref{eq:mfd-tangent-change}展开$\partial_j|_p$，得到
\[
a_j=\omega(\partial_j|_p)=\sum_k\frac{\partial\tilde x_k}{\partial x_j}\,\tilde a_k .
\]
记$\bm J$为坐标变换的Jacobian矩阵，$\bm a,\tilde{\bm a}$为两组系数构成的列向量，两种对象的分量变换规则并排写出是
\begin{equation}
\tilde{\bm v}=\bm J\bm v,\qquad \tilde{\bm a}=\bm J^{-\top}\bm a .
\label{eq:mfd-covector-change}
\end{equation}
切向量的分量随$\bm J$变换，余切向量的系数随$\bm J$的逆转置变换。二者配对的结果$\tilde{\bm a}^\top\tilde{\bm v}=\bm a^\top\bm J^{-1}\bm J\bm v=\bm a^\top\bm v$不变，这正是“余切向量是对切向量的读数”应有的性质。

\paragraph{不能未选度量就把二者等同}
在一张图中，切向量$\bm v$与系数相同的余切向量$\bm v^\top$看上去一一对应，很容易把二者当成同一个东西。但这种对应依赖图。仍取$p=(0,\tfrac35,\tfrac45)$，$\bm J=\operatorname{diag}(\tfrac19,-\tfrac19)$。北极图中的切向量$\partial_{u_1}|_p$分量为$(1,0)$，按系数相同对应到余切向量$\mathrm du_1|_p$。

现在把两者都换到南极图。由式\eqref{eq:mfd-covector-change}，切向量的分量变为$(\tfrac19,0)$，余切向量$\mathrm du_1|_p$的系数变为$\bm J^{-\top}(1,0)^\top=(9,0)$。若在南极图中重新按系数相同来对应，切向量$(\tfrac19,0)$对应的余切向量系数是$(\tfrac19,0)$，而不是$(9,0)$。两张图给出的对应不同，所以“系数相同”不是切空间与余切空间之间与图无关的对应。

要在二者之间建立与图无关的对应，需要在每个切空间上额外选定一个内积，把切向量$\bm v$对应到读数$\bm w\mapsto\langle\bm v,\bm w\rangle$。这正是黎曼度量的作用，留到第~\ref{chap:riemannian-geometry}章。本章始终把切向量与余切向量看作两种不同的对象。
